\documentclass[11pt,a4paper]{article}

\usepackage[T1]{fontenc}
\usepackage{lmodern}
\usepackage{amsmath,amssymb,amsthm,mathtools}
\usepackage{geometry}
\usepackage{graphicx}
\usepackage{booktabs}
\usepackage{caption}
\usepackage{microtype}
\usepackage{xcolor}
\usepackage{hyperref}
\usepackage{authblk}

\geometry{margin=2.3cm}
\hypersetup{
  colorlinks=true,
  linkcolor=blue!50!black,
  citecolor=green!40!black,
  urlcolor=blue!60!black,
  pdftitle={A Quantum Fourier Transform for Continuous-Variable Quantum Key Distribution},
  pdfauthor={Yong Jia and Xiaoxiang You},
}

\newtheorem{proposition}{Proposition}[section]

\newcommand{\Z}{\mathbb{Z}}
\newcommand{\eps}{\varepsilon}
\newcommand{\ii}{\mathrm{i}}

\title{\textbf{A Quantum Fourier Transform for Continuous-Variable\\
Quantum Key Distribution}}

\author[1]{Yong Jia}
\author[2]{Xiaoxiang You}

\affil[1]{School of Information Science and Technology, University of Science and Technology of China}
\affil[2]{School of Mathematical Sciences, University of Science and Technology of China}

\date{}

\begin{document}
\begin{table}[t]
\centering
\caption{The two exceptional moduli compared. $c$ is the coordination number of the deep hole, which fixes the dimension of the lowest manifold.}
\label{tab:moduli}
\small
\begin{tabular}{@{}lll@{}}
\toprule
& $\tau=\ii$ (square, $j=1728$) & $\tau=\rho$ (hexagonal, $j=0$) \\
\midrule
Unit group $\operatorname{Aut}^+(Q_\tau)$ & cyclic of order $4$ & cyclic of order $6$ \\
Relevant rotation & $S_4$, quarter turn & $S_3=S_6^2$, $120^\circ$ \\
Deep hole(s) & $(1/2,1/2)$ & $(1/3,2/3),\ (2/3,1/3)$ \\
$\delta_\tau$ & $1/2$ & $4/9$ \\
Coordination $c$ & $4$ & $3$ \\
Next shell & $5/2$ (multiplicity $8$) & $16/9$ (multiplicity $3$) \\
Shell gap & $2\kappa^2$ & $\frac{4}{3}\kappa^2$ \\
Determinant $D_\tau$ & $2$ () & $\sqrt{3}$ () \\
Susceptibility $\mathcal K$ & $2\pi I_2$ & $\frac{2\pi}{\sqrt{3}}\left(-\frac{2}{-1}-\frac{-1}{2}\right)$ \\
Natural encoding & dual-rail qubit after a selector & qutrit, no selector needed \\
\bottomrule
\end{tabular}
\end{table}

\subsection{Ideal round benchmark (square case only)}

\begin{proposition}
\label{prop:round}
The lowest eigenvalue of  is maximized uniquely at $a=(1/2,1/2)$, where for $\mu=0$ the two lowest distinct shells are $3/R^2$ (fourfold) and $8/R^2$, with separation $5/R^2$.
\end{proposition}

\begin{proof}
$t_0(a)=\operatorname{dist}(u,\Z)+\operatorname{dist}(v,\Z)\leq 1$ with equality only at the antiperiodic point, and $\ell\mapsto \ell(\ell+2)$ is increasing.
\end{proof}

This model minimizes an $\ell^1$ quantity, not a quadratic form, so neither nor  is used: the optimum is the elementary maximization of $\operatorname{dist}(u,\Z)+\operatorname{dist}(v,\Z)$. It has no hexagonal analogue. Its gaps are properties of the Friedrichs extension at the two singular axes; other self-adjoint extensions exist, and a finite core interpolates, so the numbers are idealized benchmarks.

\section{Measurable predictions}
\label{sec:predictions}

The determinant of  is a collective one-loop quantity and is not a laboratory observable; we therefore do not propose it as a test. The measurable content is the resonance-frequency pattern.

\begin{proposition}[Linear splitting, stationary centroid]
\label{prop:splitting}
Let $a_0$ be a deep hole of $Q_\tau$ with nearest offsets $w_i=k_i+a_0$, $i=1,\ldots,c$, and let $a=a_0+\eps$. Then, exactly,
\begin{equation}
Q_\tau(w_i+\eps)=\delta_\tau+2\,w_i^TG_\tau^{-1}\eps+Q_\tau(\eps),
\label{eq:splitting}
\end{equation}
so the manifold splits linearly. If, in addition, $\sum_i w_i=0$, as at the square and hexagonal symmetry-fixed deep holes, the centroid obeys
\begin{equation}
\overline Q=\delta_\tau+Q_\tau(\eps),\qquad \operatorname{Hess}\,\overline Q=2\,G_\tau^{-1}.
\label{eq:centroid}
\end{equation}
These branches are the $c$ lowest eigenvalues provided $\|\eps\|_{Q_\tau}<\tfrac12(\sqrt{\delta^{(2)}}-\sqrt{\delta_\tau})$, where $\delta^{(2)}$ is the next shell value; for $\tau=\ii$ the sharp condition is
\begin{equation}
2|\eps_u|+|\eps_v|<1\quad\text{and}\quad |\eps_u|+2|\eps_v|<1,
\label{eq:sharp}
\end{equation}
for which $Q=\tfrac12\pm\eps_u\pm\eps_v+\eps_u^2+\eps_v^2$ with independent signs. In the round model the four lowest eigenvalues are $(1+s)(3+s)/R^2$ with $s=\pm\eps_u\pm\eps_v$, and $\tilde\lambda=(3+\eps_u^2+\eps_v^2)/R^2$, valid precisely for $|\eps_u|+|\eps_v|<1/2$.
\end{proposition}
\end{document}
